\chapter{Validation and Its Limits}
\label{chap:validation}

A simulator that agrees with itself is not validated. This chapter states
precisely what has been established, what has been retracted, and what currently
rests on nothing external. It is deliberately the least flattering chapter in
this thesis, because the alternative --- presenting internal consistency as
physical evidence --- is how simulation results become unfalsifiable.

\section{A taxonomy of checks}

It is useful to separate four kinds of evidence, in increasing order of strength:

\begin{center}
\begin{tabular}{@{}l p{0.62\textwidth}@{}}
\toprule
\textbf{Kind} & \textbf{What it establishes} \\
\midrule
Unit / numerical
& That a routine computes the formula it claims to. Necessary, never sufficient.
  Includes machine-precision identities and mutation-tested gates. \\
\addlinespace
Internal consistency
& That two independent code paths agree where they must --- e.g.\ the multitone
  solve reproducing the Floquet solve in the small-signal limit. Strong evidence
  \emph{about the implementation}, silent about the model. \\
\addlinespace
Cross-code agreement
& That this solver and another solver produce the same number. Establishes that
  two implementations of the same model agree. \textbf{Not} physical evidence. \\
\addlinespace
External physical
& That the model reproduces a measurement of a real device, with independently
  calibrated inputs. The only kind that validates physics. \\
\bottomrule
\end{tabular}
\end{center}

Most of what exists for this solver is in the first two rows.

\section{What is established}

\subsection{Numerical identities}

Several quantities were verified to machine precision against independently
constructed references, as reported in \cref{chap:pumphb}:
the dynamic block $\Dmat_k$ (relative deviation $\num{0}$); the tangent-stiffness
harmonic $\Khat_\ell$ ($\num{0}$); the convergence norm (twelve digits); the
DC-subtracted branch current ($\num{2.6e-23}$).

The fast preconditioner backend reproduces the reference assembly to
$\num{1.0e-11}$ relative, tested quadrant by quadrant --- necessary because a
sign error there would only make the preconditioner slower, and every physics
gate would still pass.

\subsection{Conjugate symmetry}

$\ghat_{-\ell} = \overline{\ghat_\ell}$ is an exact identity for a real pump
waveform (\cref{eq:conj-symmetry}), reported per sideband in
\code{gamma\_hat\_summary.csv}. It is the sharpest available check that the
pump waveform, its basis metadata, and its phase convention all agree.

\subsection{Component reproducibility}

The production circuit rebuilds bit-identically from its own stored parameters:
$16312/16312$ elements matching, with maximum matrix difference exactly $0.0$.
The scatter RNG stream is pinned so that historical results reproduce at the
same seed.

\subsection{Basis-convergence, partially}

\begin{measured}[JTWPA lattice family, 8 of 9 settings]
\begin{center}
\begin{tabular}{@{}l l r l@{}}
\toprule
Knob & Top two settings & $|\Delta P_{1\mathrm{dB}}|$ & Verdict \\
\midrule
Signal order $Q$ & $Q{=}2$, $Q{=}3$ & \SI{0.059163}{\dB} & PASS \\
Torus scale & $\text{sc}{=}1$, $\text{sc}{=}2$ & \SI{0.000000}{\dB} & PASS \\
Pump order & order 1, order 3 & \SI{1.656572}{\dB} & \textbf{not evaluable} \\
Odd vs dense & dense order-5 did not run & --- & \textbf{not evaluated} \\
\bottomrule
\end{tabular}
\end{center}
\end{measured}

Torus scale is converged \emph{exactly}: doubling $N_p$ and $N_\delta$ moves
nothing, confirming the alias-free grid sizing of \cref{chap:saturation}.

The pump-order number spans a gap and is therefore not a convergence
measurement: the odd order-2 basis $\{1,3\}$ fails its pump solve reproducibly
in $\approx\SI{1.4}{\second}$ across two independent runs, while $\{1\}$ and
$\{1,3,5\}$ both converge. Orders 1 and 3 are not adjacent settings, so
\SI{1.66}{\dB} between them measures nothing. The dense order-5 point --- the
mandatory odd-versus-dense comparison --- exhausted memory.

A consistency check that did pass: \code{three\_tone} and \code{lattice} at
$Q{=}1$, order 1 agree to all six reported decimals ($-112.748281$), as they must,
being literally the same basis.

\begin{caveat}[this does not discharge the production-basis question]
Every setting in that matrix solves at $3.756$--\SI{7.894}{\dB} small-signal
gain, against production's \SI{27.541}{\dB}, and the study builds
\code{build\_lattice\_basis} where production uses
\code{build\_sideband\_matched\_basis} at $S=10$. It converges the lattice family
at a much weaker operating point. The production basis has \emph{not} passed a
\SI{0.2}{\dB} convergence gate and must not be described as having done so.
\end{caveat}

\section{The domain of validity of the harmonic-balance ansatz}
\label{sec:ansatz-validity}

Every check so far asks whether the solver solves its equations correctly. This
section asks a prior question: over what range of operating points is the
harmonic-balance \emph{ansatz} --- amplitudes on a fixed frequency lattice ---
capable of representing the waveform at all. No amount of Newton convergence
answers that, because a converged solution on an inadequate lattice is
converged to the wrong problem.

\subsection{The instrument}

The measurement needs a reference that assumes no lattice. The time-domain
kernel of \cref{chap:numerics} is such a reference: it integrates the node
equations directly, so its spectrum is whatever the circuit produces. Given a
steady-state trace, define

\begin{definition}[on-lattice fraction]
For a pumped, signal-driven circuit with tone lattice
$\omega_{(h,q)} = h\wp + q\ws$ (\cref{eq:tone-frequency}), the \emph{on-lattice
fraction} is the share of in-band spectral power lying within one analysis
window of some lattice node:
\begin{equation}
  \Lambda = \frac{\sum_{\,\omega \,\in\, \mathcal{W}} P(\omega)}
                 {\sum_{\omega} P(\omega)},
  \qquad
  \mathcal{W} = \Bigl\{\omega : \min_{h,q}\,
                 \bigl|\omega - \omega_{(h,q)}\bigr| \le w \Bigr\}.
  \label{eq:on-lattice-fraction}
\end{equation}
\end{definition}

$\Lambda$ is an upper bound on any harmonic-balance result at that operating
point, and it requires no calibration: it is a ratio of powers within one
spectrum, so the port convention, the line-loss model and the absolute power
scale all cancel.

Two refinements separate the ways the bound can fail. Admitting half-integer and
third-integer pump nodes gives $\Lambda_{1/2}$ and $\Lambda_{1/3}$: a
period-doubled state would lift $\Lambda_{1/2}$ close to unity, so if it does
not, a period-2 basis cannot recover the missing power. And the residue
$1-\Lambda$ is characterised by the share of off-lattice power held by its
twenty strongest bins, and by a scan over one extra generator $f_a$ asking how
much of the residue falls on $n\wp + m\ws + k\,\omega_a$. Discrete residue
concentrates and fits a generator --- the signature of a torus, representable by
a three-frequency balance. A raised continuum does neither and is representable
by nothing.

The measurement is \code{scripts/chaos/measure\_ansatz\_validity.py}; the
campaign it reduces covers four devices and \num{68} operating points.

\subsection{Result: the ansatz is exact where the device is useful}

\begin{measured}[on-lattice fraction below the transition]
\begin{center}
\begin{tabular}{@{}l l r r@{}}
\toprule
Device & Control range & $\Lambda$ & Gain [\si{\dB}] \\
\midrule
\code{jc\_jtwpa}     & $-30.5 \to \SI{-28.2}{dBm}$ & $1.0000 \to 0.9233$ & $13.9 \to 36.1$ \\
\code{jc\_fqjtwpa}   & $-33.7 \to \SI{-31.5}{dBm}$ & $1.0000 \to 0.9346$ & $27.1 \to 31.0$ \\
\code{ipm\_2c\_fixed}& $0.450 \to 0.575\,I_{\mathrm{bound}}$ & $1.0000$ & $8.2 \to 13.2$ \\
\code{guarcello}     & $-70 \to \SI{-54}{dBm}$ & $1.0000$ & $0.02 \to 8.5$ \\
\bottomrule
\end{tabular}
\end{center}
The off-lattice floor over these ranges sits \num{112} to \SI{226}{\dB} below the
pump.
\end{measured}

Across every device, throughout the amplifying regime, the waveform is on the
production tone lattice to four decimal places. Whatever else limits the results
of \cref{chap:saturation}, the choice of ansatz does not.

\subsection{Result: the transition is not a period doubling}

Above a sharply defined threshold the two travelling-wave devices lose the
lattice entirely, within a single \SI{0.4}{dB} step of pump power.

\begin{measured}[at and above the transition]
\begin{center}
\begin{tabular}{@{}l r r r r@{}}
\toprule
Lattice admitted & \code{jc\_jtwpa} & \code{jc\_fqjtwpa} & \code{ipm\_2c\_fixed} & \code{guarcello} \\
\midrule
$n\wp + m\ws$              & $0.056$ & $0.051$ & $0.76$--$0.91$ & $0.90$--$0.98$ \\
$+$ half-integers          & $0.104$ & $0.095$ & $+0.005$ & $+0.004$ \\
$+$ thirds                 & $0.151$ & $0.159$ & $+0.008$ & $+0.006$ \\
\midrule
Top-20 share of residue    & $0.024$ & $0.032$ & $0.36$--$0.72$ & $0.27$ \\
Best single generator      & $0.262$ & $0.268$ & $0.52$--$0.78$ & $0.475$ \\
\bottomrule
\end{tabular}
\end{center}
\end{measured}

Admitting half-integer pump nodes recovers \num{4.8} and \num{4.5} percentage
points on the two travelling-wave devices, and under one point on the other two.
A period-doubled state would have produced the opposite result.

\begin{caveat}[no period-\emph{N} ansatz is justified by this evidence]
The dormant period-doubling scaffolding of \cref{chap:saturation} ---
\code{build\_half\_pump\_basis}, the period-doubled pump basis, its branch
continuation and its small-signal solve --- stays dormant. The measurement above
is the direct test of the hypothesis those modules exist to serve, and it fails
it. Enabling them would require a tracked Floquet multiplier crossing, not a
loss of convergence.
\end{caveat}

Nor is the collapsed state a torus. Its twenty strongest off-lattice bins hold
\SI{2.4}{\percent} of the off-lattice power, and the best single extra generator
explains \SI{26}{\percent} --- both consistent with a continuum, whose floor sits
only $13$ to \SI{18}{\dB} below the pump. This is the regime in which harmonic
balance of any order has nothing to offer, and the time-domain kernel is the
only available instrument.

\subsection{Result: a torus exists, and the window is narrow}

Between the two regimes there is a band in which the residue is small and
\emph{discrete}.

\begin{measured}[\code{jc\_jtwpa}, approaching the transition]
\begin{center}
\begin{tabular}{@{}r r r r r@{}}
\toprule
$P_p$ [\si{dBm}] & Gain [\si{\dB}] & $1-\Lambda$ & Top-20 share & Generator fit \\
\midrule
$-29.3$ & $27.4$ & $0.009$ & $0.948$ & $0.999$ \\
$-28.9$ & $31.4$ & $0.027$ & $0.869$ & $0.999$ \\
$-28.6$ & $34.1$ & $0.044$ & $0.612$ & $0.872$ \\
$-28.2$ & $36.1$ & $0.077$ & $0.371$ & $0.627$ \\
$-27.8$ & $-6.0$ & $0.944$ & $0.024$ & $0.262$ \\
\bottomrule
\end{tabular}
\end{center}
\end{measured}

A second incommensurate frequency is born roughly \SI{1.1}{\dB} of pump power
before the collapse, and never carries more than \SI{7.7}{\percent} of the power
before the continuum replaces it. A quasi-periodic formulation --- the existing
two-frequency lattice with the second frequency promoted from a prescribed
detuning to an unknown, closed by an auxiliary-generator condition --- would
extend the representable range by about that much and no further. That is a real
extension with a measured, and modest, size.

\subsection{An unresolved gain disagreement, and how it was narrowed}
\label{sec:gain-gap}

The same comparison that validates the ansatz exposes a disagreement in the
quantity the solver exists to compute. At matched operating points, harmonic
balance and the time-domain kernel do not report the same gain.

The raw discrepancy is large, and four candidate explanations were tested. Three
are settled.

\paragraph{The analysis window is not the cause.} The campaign fits the signal
tone over the last \num{100} pump periods. Re-extracting the identical traces
over \num{200}, \num{400} and \num{800} periods moves the result by at most
\SI{0.5}{\dB} throughout the amplifying regime.

\paragraph{Sideband truncation is not the cause.} At the matched point, the
multitone small-signal gain is converged in $S$: \SI{0.012}{\dB} between $S=10$
and $S=12$ (\cref{sec:jc-retraction}).

\paragraph{The probe was partly compressing, and the kernel has a signal-tone
floor.} The campaign drove a probe of \SI{3e-8}{\ampere} against a pump of
\SI{3.9e-6}{\ampere}. Lowering the probe by successive decades reveals both a
real compression and a hard measurement floor:

\begin{measured}[\code{jc\_jtwpa}, $f_p=\SI{7.12}{\giga\hertz}$: output amplitude at $\ws$]
\begin{center}
\begin{tabular}{@{}r r r r l@{}}
\toprule
Probe [\si{\ampere}] & Drive [\si{\volt}] & \multicolumn{1}{c}{$\SI{-30.1}{dBm}$} & \multicolumn{1}{c}{$\SI{-29.7}{dBm}$} & Behaviour \\
\midrule
\num{3e-8}  & \num{1.50e-6} & \num{6.07e-6} & \num{1.10e-5} & signal \\
\num{3e-9}  & \num{1.50e-7} & \num{7.46e-7} & \num{1.47e-6} & signal $+$ floor \\
\num{3e-10} & \num{1.50e-8} & \num{2.73e-7} & \num{3.52e-7} & floor-dominated \\
\num{3e-11} & \num{1.50e-9} & \num{2.34e-7} & \num{2.57e-7} & floor \\
\num{3e-12} & \num{1.50e-10} & \num{2.30e-7} & \num{2.49e-7} & floor \\
\bottomrule
\end{tabular}
\end{center}
Over four decades of drive the output stops responding below
\SI{3e-10}{\ampere} and settles at \num{2.3e-7} to \SI{2.5e-7}{\volt}, moving
\SI{1.6}{} to \SI{3.4}{\percent} across the last factor of ten. The
\SI{3e-12}{\ampere} probe accordingly reports a ``gain'' of \SI{63.7}{} and
\SI{64.4}{\dB}. That level does \emph{not} fall as the analysis window
lengthens, so it is not window leakage: it is undamped content at the signal
frequency, consistent with the ramp transient that this kernel's
non-\hbox{$L$-stable} integrator never removes (\cref{chap:numerics}). It is
slightly higher at the higher pump power, as parametrically amplified residue
should be.
\end{measured}

Two consequences. Probes at or below \SI{3e-10}{\ampere} measure the floor and
their gains --- which rise without bound as the probe shrinks --- are spurious.
And after subtracting the floor in quadrature, the \SI{3e-8}{\ampere} probe is
compressed by \SI{1.2}{\dB} at \SI{-30.1}{dBm} and \SI{2.3}{\dB} at
\SI{-29.7}{dBm}.

\paragraph{The reference column was not converged in $S$.} The harmonic-balance
column this comparison was first made against inherits \code{--sidebands 6}, the
gain-map default, which nothing in that driver overrides. Its numbers are
therefore $S=6$, roughly \SI{1.4}{\dB} below the converged value. Re-solving at
the \emph{same} pump current confirms both facts at once:

\begin{measured}[multitone versus Floquet at $I_p = \SI{3.8806468570637416e-6}{\ampere}$]
\begin{center}
\begin{tabular}{@{}l r@{}}
\toprule
Path & Gain [\si{\dB}] \\
\midrule
Floquet column, $S=6$      & $29.140947$ \\
Multitone, $S=6$           & $29.141$ \\
Multitone, $S=10$          & $30.542$ \\
Multitone, $S=12$          & $30.554$ \\
\bottomrule
\end{tabular}
\end{center}
Same-$S$ parity holds to \SI{0.0002}{\dB}, and the pump utilisation agrees to
eleven significant figures ($0.6167509694508315$ against $0.6167509694519268$).
\end{measured}

Same-$S$ multitone/Floquet parity is therefore intact, as the unit gate at
$S=10$ on the JPA fixture already asserted to \SI{1e-4}{\dB}. What the reference
column measures is a $S=6$ truncation, and comparing it against a converged
multitone solve is not a parity test.

Correcting for that, and comparing at the time-domain run's \emph{own} achieved
on-chip pump current so that no power convention enters:

\paragraph{The cause was the kernel's timestep.} The one candidate that could
not be excluded by inspection was the time-domain step itself, which for this
device has \emph{no} independent validation: the linear-limit check of
\cref{chap:numerics} is degenerate on \code{jc\_jtwpa}, because setting $I_c=0$
removes its only inductive path, so both sides return exactly zero. Halving the
step settles it.

\begin{measured}[gain versus timestep, at a \SI{3e-9}{\ampere} probe]
\begin{center}
\begin{tabular}{@{}r r r r r r@{}}
\toprule
$P_p$ [\si{dBm}] & $I_p$ [\si{\ampere}] & $\Delta\theta=0.01$ & $\Delta\theta=0.005$ & Model $S{=}12$ & Residual \\
\midrule
$-30.1$ & \num{3.701475e-6} & $19.996$ & $24.916$ & $27.743$ & \SI{2.83}{\dB} \\
$-29.7$ & \num{3.873843e-6} & $25.891$ & $29.585$ & $30.440$ & \SI{0.86}{\dB} \\
\bottomrule
\end{tabular}
\end{center}
All gains are \code{gain\_vs\_off\_db}, a pump-on/pump-off ratio, and are
therefore invariant under the port-power convention and the line-loss model.
Halving the step moves the time domain \SI{+4.9}{} and \SI{+3.7}{\dB}, toward
the model in both cases, cutting the disagreement by factors of \num{2.7} and
\num{5.3}. Junction utilisation moves by less than \SI{0.5}{\percent} over the
same refinement ($0.5856 \to 0.5883$ and $0.6114 \to 0.6125$).
\end{measured}

That last line is the diagnostic point. The pump state was never in question ---
it agreed with harmonic balance to better than \SI{1}{\percent} at both steps.
What was under-resolved is the small-signal response, which is exactly where the
utilisation agreement had already localised the fault.

\begin{caveat}[the timestep is not yet converged]
The residual disagreement is still falling with the step and the sequence has
only two points, so no convergence order can be fitted and the remaining
\SI{0.9}{}--\SI{2.8}{\dB} must not be reported as a physical discrepancy. The
precedent on the one device where the linear-limit check does apply is that the
required step was four times finer than the default; a third point at
$\Delta\theta = 0.0025$ is required before any residual gap is claimed.

Two corrections follow for anyone reading earlier drafts. A \SI{4.8}{}--%
\SI{8.4}{\dB} gap ``with the model reading high'' was reported here on the
$\Delta\theta = 0.01$ data and is withdrawn: most of it was discretisation. And
the floor-corrected time-domain values that gap was computed from carried a
\SI{\pm1.4}{\dB} quadrature assumption that is now moot, since the timestep
effect is several times larger than the correction it was competing with.
\end{caveat}

A separate inconsistency surfaced by the same comparison: the recorded pump
utilisation for FQJTWPA runs \num{0.086}--\num{0.151} across its whole converged
column while the time domain reports \num{0.62}--\num{0.89} at the same points,
and the harmonic-balance sequence is not even monotone in pump power. The JTWPA
column, checked above, shows no such problem. The FQJTWPA figure is not usable
as a boundary diagnostic until that is explained.

\subsection{What this section does and does not establish}

\begin{caveat}[an independent implementation, not an independent measurement]
The time-domain kernel and the harmonic-balance solver share this repository's
circuit matrices, builders and loss models. They do not share a frequency
ansatz, a time discretisation or a solution method, so the comparison is
informative about the ansatz specifically. It remains cross-code agreement in
the sense of \cref{sec:jc-retraction} and is \emph{not} external physical
validation.
\end{caveat}

The kernel's own timestep is verified per device against an analytic linear
limit (\cref{chap:numerics}); \code{ipm\_2c\_fixed} passes at
\SI{0.76}{\percent}. That check is degenerate on the two \code{jc\_*} devices,
which therefore carry no independent linear validation of their own --- a gap
this section inherits rather than closes.

\section{What has been retracted}

\subsection{Cross-code agreement is not validation}
\label{sec:jc-retraction}

The sideband truncations used in production ($S=10$ for JTWPA, $S=6$ for
FQJTWPA, $S=10$ for 2c) were originally selected by gating small-signal gain
against a reference implementation in another language, to \SI{0.2}{\dB}.

That gate is void, for two independent reasons:

\begin{enumerate}
\item The reference is another simulator. It has no reference of its own. Two
      codes agreeing establishes that they implement the same equations, which
      is a useful \emph{regression} check and nothing more.
\item Two of the seven parity designs are that code's own documentation
      examples. Validating a solver against designs distributed by the code it
      is being compared to is circular.
\end{enumerate}

\begin{caveat}[what this leaves unresolved]
JTWPA gain is \textbf{non-monotone in $S$}:
\begin{center}
\begin{tabular}{@{}r r r r r@{}}
\toprule
$S$ & 2 & 4 & 6 & 10 \\
\midrule
Gain [\si{\dB}] & $30.7152$ & $24.2021$ & $26.5563$ & $27.5410$ \\
\bottomrule
\end{tabular}
\end{center}
Without an external reference, $S=10$ cannot be selected from that sequence by
agreement. Selecting it requires a self-convergence argument. Non-monotonicity
means $S=10$ may still be climbing. The 2c device never had a cross-code
reference at all and inherited $S=10$ by analogy.
\end{caveat}

That argument has since been run for both devices, and it succeeds. For 2c, at
$f_p = \SI{7.100}{\giga\hertz}$ and three signal frequencies, $S=10$, $12$ and
$14$ agree on both $G_0$ and $P_{1\mathrm{dB}}$ to below \SI{1e-6}{\dB}. For
JTWPA, at $f_p = \SI{7.12}{\giga\hertz}$ and
$I_p = \SI{3.873843e-6}{\ampere}$:

\begin{measured}[JTWPA small-signal gain versus sideband count]
\begin{center}
\begin{tabular}{@{}r r r r r r@{}}
\toprule
$S$ & 4 & 6 & 8 & 10 & 12 \\
\midrule
Gain [\si{\dB}] & $27.539$ & $29.067$ & $30.300$ & $30.428$ & $30.440$ \\
\bottomrule
\end{tabular}
\end{center}
Monotone throughout, with $S=10 \to S=12$ moving \SI{0.012}{\dB}.
\end{measured}

\begin{caveat}[what the sequence above does and does not settle]
It is measured at one operating point on one device, with the pump basis held
at its production value. It establishes that $S=10$ is converged \emph{there},
which is what the results of \cref{chap:saturation} require. It does not explain
the non-monotone sequence quoted above, which was measured at a different
operating point and remains unexplained.

Same-$S$ multitone/Floquet parity remains a valid internal check and should
still pass.
\end{caveat}

The parity results themselves are retained as \emph{provenance}: six of seven
designs match to $<\SI{0.0024}{\dB}$, and one lossy design remains
$\approx\SI{0.89}{\dB}$ off pending reconciliation of the lossy pump convention.
That is a statement about numerical drift between two codes. It is reported as
such.

\subsection{The stability result}

The dominant-exponent result reported in an earlier draft is withdrawn; see
\cref{chap:saturation}. Two defects, both now fixed and mutation-verified.
The instructive part is \emph{why} the tests missed them: they ran on $I_c=0$,
which is a linear circuit, at a zero state. A test that exercises no
nonlinearity cannot detect an error in the nonlinear tangent.

\subsection{\texorpdfstring{The $P_{1\mathrm{dB}}$ refinement discrepancy}
                          {The P1dB refinement discrepancy}}

An earlier claim of $+0.461/+\SI{2.920}{\dB}$ between interpolated and refined
$P_{1\mathrm{dB}}$ is withdrawn: it was measured on a coarse grid no published
number used, which placed one crossing \SI{8.4}{\dB} off. The correct figures
are $0.23$--$\SI{0.34}{\dB}$ (\cref{chap:saturation}).

\section{What has no external reference}

\subsection{Saturation}

\begin{caveat}[saturation currently rests on internal checks only]
Two candidate external references were considered and both are unavailable:
\begin{itemize}
\item The cross-code reference is retired for the reasons in
      \cref{sec:jc-retraction}, and in any case is regression-only.
\item The available measurement cube from the cryogenic setup sweeps pump power
      against signal frequency. It has \textbf{no signal-power axis}, so it
      cannot validate saturation or $P_{1\mathrm{dB}}$ at all. (It remains a
      valid reference for gain maps.)
\end{itemize}

Saturation correctness therefore rests entirely on:
\begin{enumerate}
\item the conversion-scope Manley--Rowe invariant (valid, with a real
      $\approx\SI{2.2}{\percent}$ floor);
\item power balance;
\item production-basis self-convergence (\textbf{never measured});
\item the small-signal Floquet limit (a genuine cross-path check).
\end{enumerate}
Items 3 and 4 are the ones that would carry weight, and item 3 has not been done.
\end{caveat}

A later dataset does sweep signal power and would constitute the first genuine
external compression reference. Reconciling its calibration and gain-reference
conventions is future work.

\subsection{The truncation bias}

\begin{measured}[the matched basis caps $|q|$ at 1]
\code{build\_sideband\_matched\_basis} produces tones with $|q|\le1$ at
\emph{every} sideband count. Opening $q$ to $\pm2$ cuts compression by
\SI{44}{\percent}.
\end{measured}

This is a definite, signed bias: the published $P_{1\mathrm{dB}}$ values are
\emph{low}. It is a property of the basis constructor, not of the solver, and it
is the most concrete known limitation of the compression results.

\section{Comparison with measurement}

\subsection{Fitting calibration offsets}

Comparing a simulated map to a measured one requires aligning calibration
offsets in frequency and power. Rather than tune them by hand, they are fitted
as nuisance parameters:
\begin{equation}
  G_{\mathrm{meas}}(f,P) \approx
  G_{\mathrm{sim}}(f - \Delta f,\; P - \Delta P) + \Delta G .
  \label{eq:map-alignment}
\end{equation}
For weighted least squares $\Delta G$ is analytic given $(\Delta f,\Delta P)$, so
only a two-dimensional grid search remains. Non-overlapping and failed cells are
masked; the amplified ridge is weighted above the flat background so that the
fit is driven by the feature of interest.

\begin{measured}[per-section fits are far better identified]
A whole-map fit gives $\Delta f = \SI{-0.30}{\giga\hertz}$,
$\Delta P = \SI{+2.55}{\dB}$, $\Delta G = \SI{-1.30}{\dB}$, RMS \SI{6.0}{\dB},
with a frequency-elongated (weakly identified) basin --- because a single
$\Delta f$ cannot align all comb lobes when the comb phase drifts.

Restricting to one comb branch gives $\Delta f \approx 0$ --- the two datasets'
combs are already frequency-aligned --- and
$\Delta P \approx +2.5$ to $\SI{+3.3}{\dB}$ robustly across
$6.2$--\SI{7.45}{\giga\hertz}, with a compact single minimum. That is the one
real calibration offset.
\end{measured}

Per-lobe RMS remains $2$--\SI{4}{\dB}, for two reasons that are model-fidelity
and coverage gaps rather than fit failures: the simulation does not reach the
measured high-power lobes (a numerical-boundary cap), and its comb periodicity
differs slightly.

\begin{implnote}[guard against tiny-overlap fits]
A soft $1/\text{overlap}$ penalty alone lets a $\approx9$-cell corner with
near-zero local residual win the fit. \code{--min-overlap-frac} (default $0.25$)
rejects those outright. The threshold is relative to the floor-weighted window,
whose maximum achievable overlap is only $0.3$--$0.4$ because the simulation
fails at high power --- so $0.25$ is the practical ceiling, not a lax setting.
\end{implnote}

\subsection{The model-fidelity gap}

\begin{measured}[gain-map reproduction]
Power-substep recovery reduced the failed-cell count on a measurement-matched
map from 1017 to 103 --- i.e.\ the \emph{numerical} boundary was largely solved.
The model nonetheless undershoots the measured peak gain by
$\approx\SI{6.5}{\dB}$. That residual is a model-fidelity gap, not a
convergence problem.
\end{measured}

\begin{measured}[compression, six frequencies]
Against the live device, the model compresses \emph{late} at all six measured
frequencies, by $+1.15$ to $+\SI{5.51}{\dB}$. The gain band itself matches well
(a uniform \SI{0.92}{\dB} deficit). Ripple makes few-frequency comparisons
unreliable: the spread across frequencies exceeds the mean offset.
\end{measured}

\begin{caveat}[compare at equal gain, and quote a slope]
A single-frequency comparison of $P_{1\mathrm{dB}}$ between model and
measurement is nearly meaningless, because $P_{1\mathrm{dB}}$ depends strongly on
gain and both datasets sweep frequency, not gain. The defensible comparison is
the \emph{slope} of $P_{1\mathrm{dB}}$ against gain.

Even that is confounded: every available slope on both sides is a frequency
sweep in disguise. The model's slopes span $-0.63$ to $-3.85$; at $G>\SI{25}{\dB}$
the JTWPA slope is $-3.85$, steeper than the hardware's. A measured
$\approx\SI{2.6}{\dB\per\dB}$ gap survives testing against estimator artefacts,
shared gain-reference noise, and five other candidate mechanisms.
\end{caveat}

\subsection{Pump calibration}

\begin{measured}[on-chip pump confirmed]
The on-chip pump power is confirmed to \SI{0.7}{\dB} by two independent routes
($\SI{-65.98}{\dBm}$ against $\SI{-66.7}{\dBm}$). The apparent
$\approx\SI{7.9}{\dB}$ gap between model and measurement is therefore not a
calibration error.
\end{measured}

Of that gap, $\SI{6.02}{\dB}$ is the Norton-termination factor of
\cref{chap:circuit}: the ports are Norton-terminated, so computing on-chip power
as $I^{2}Z_0/2$ overstates it. Correcting for it leaves $\approx\SI{1.1}{\dB}$.

\begin{caveat}[line loss must be applied]
Two campaigns ran with \emph{no} pump-line loss at all. The measured loss model
of \cref{chap:circuit} gives \SI{34.54}{\dB} at \SI{7.1}{\giga\hertz}; omitting
it is not a small correction. All future work uses the measured model.
\end{caveat}

\begin{implnote}[on-chip dielectric-loss gate]
With the pump off, the on-chip tangent is checked independently of another
simulator: $\ln|S_{21}|$ must scale linearly with $\tan\delta$ and
$|S_{21}|$ must decrease monotonically. This is the analytic
$\exp(-\alpha L)$ attenuation check, separate from measured input-line loss.
It does not make saturation an externally validated result.
\end{implnote}

\section{Two structural findings}

\subsection{The device is not travelling-wave}

\begin{measured}
For the two-chip device, $|V^-|/|V^+| = 0.359$: it is a standing-wave resonant
structure. This is \emph{not} a port mismatch --- two earlier explanations,
based on the midpoint impedance and on an \SI{84.6}{\ohm} characteristic
impedance, were both checked and both wrong.
\end{measured}

The consequence is interpretive rather than numerical: the solver computes the
full scattering solution correctly either way, but coupled-mode expressions of
the form \eqref{eq:gain-cme}, and any reasoning that assumes a monotonically
growing envelope, do not apply.

\subsection{Design directories are the source of truth}

\begin{caveat}[an entire campaign ran on a stale circuit]
Output directories from previous runs contain circuit matrices, and they are
easy to mistake for designs. They are not: they are snapshots. A series of
compression campaigns ran on a stale circuit taken from an output directory
rather than from the live design.

The live designs are under \file{designs/}. This is why the design directory
convention, the topology gate of \cref{sec:profiles}, and the rebuild-from-stored
parameters property all matter operationally, not just aesthetically.
\end{caveat}

\section{Testing methodology}

\subsection{Mutation testing}

Every gate in the component-profile and scatter suite (95 tests) was verified by
\emph{mutation}: deliberately breaking the implementation and confirming the test
fails. A test that has never been seen to fail is not evidence.

This methodology directly caught the two stability defects described above ---
the original tests passed on a linear circuit at a zero state and could not have
detected either.

\subsection{Clone testing}

Local test runs can pass spuriously because files remain on disk after being
removed from version control. A dependency dropped from the repository is still
importable locally. Verification therefore includes a clean-clone run: check out
the repository fresh, install, and run the suite.

This caught a missing package that every local run had silently satisfied.

\subsection{Verifying agent-assisted work}

Where automated tooling reports work as complete, the claim is checked against
the object database (\code{git ls-tree}) rather than accepted from a summary. On
one occasion the repository head was unimportable for eight consecutive commits
despite each being reported as complete and tested.

\section{Summary of the validation position}

\begin{center}
\begin{tabular}{@{}l l p{0.42\textwidth}@{}}
\toprule
\textbf{Quantity} & \textbf{Status} & \textbf{Basis} \\
\midrule
Pump periodic steady state
& Established
& Machine-precision identities; conjugate symmetry; residual convergence. \\
\addlinespace
Passive S-parameters
& Established
& Direct linear solve, explicit residual check. \\
\addlinespace
Small-signal gain (given a basis)
& Established numerically
& Explicit linear residual; internal consistency. Cross-code agreement
  reported as regression only. \\
\addlinespace
Choice of sideband count $S$
& \textbf{Open}
& Selection gate retired; self-convergence never measured; gain non-monotone
  in $S$. \\
\addlinespace
Gain map versus measurement
& Partial
& Numerical boundary largely resolved; $\approx\SI{6.5}{\dB}$ peak-gain
  model-fidelity gap remains. \\
\addlinespace
$P_{1\mathrm{dB}}$ / compression
& \textbf{Open}
& No external reference available; known signed low bias from the $|q|\le1$
  truncation; basis convergence unmeasured. \\
\addlinespace
Large-signal stability
& Preliminary
& Earlier result withdrawn; current implementation fixed and
  mutation-verified; no unstable point found. \\
\bottomrule
\end{tabular}
\end{center}

The honest summary is that the \emph{numerics} are in good shape and the
\emph{physics validation} is incomplete, with the incompleteness concentrated
exactly where the scientific interest is highest --- the saturation regime.
Stating that plainly is more useful than a validation section that reads well.
